\section{Proof schedules and graph compilation}\label{sec:graph}

\subsection{Local rules versus global proofs}

A normalized graph node does not merely compute a value.  It exports a local rule
\[
  \Delta_v\vdash[\PRAAx]\phi_v,
\]
where \(\Delta_v\) is the finite list of predecessor certificates actually used by the node.  A topological schedule turns these local rules into a global proof.  The schedule is a metalevel object, so its dependency lists may be ordinary Lean finite lists and its acyclicity may be checked by the elaborator.

\begin{definition}[Kernel proof schedule]
Fix a context \(\Gamma\).  A proof schedule for \(\varphi\) is generated by:
\begin{enumerate}[label=(\roman*)]
  \item a terminal derivation \(\Gamma\vdash\varphi\); or
  \item a derivation \(h:\Gamma\vdash\psi\), followed by a schedule for \(\varphi\) under the enlarged context \(\psi::\Gamma\).
\end{enumerate}
The new formula \(\psi\) may be an equation defining a register, a reflected Boolean fact, a loop invariant instance, or a validity certificate.
\end{definition}

\begin{theorem}[Schedule compilation]\label{thm:schedule}
Every kernel proof schedule compiles to a derivation \(\Gamma\vdash\varphi\).
\end{theorem}

\begin{proof}
Induct on the schedule.  A terminal schedule already contains the desired derivation.  In the extension case, the induction hypothesis gives \(\psi::\Gamma\vdash\varphi\).  Apply simultaneous cut: discharge the first hypothesis with \(h\), and discharge each inherited member of \(\Gamma\) by the corresponding hypothesis rule.  The result is \(\Gamma\vdash\varphi\).
\end{proof}

This tiny theorem is the object-proof core of the graph compiler.  Everything else constructs the node formulas and their local derivations.

\subsection{Static graph elaboration}

A \KIFC graph carries the following finite metadata:
\begin{itemize}
  \item node identifiers in topological order;
  \item a node class and theorem-specific parameters;
  \item predecessor identifiers;
  \item a demand transformer for each incoming observation;
  \item register widths and state projections;
  \item the local formula exported by the node;
  \item a Lean expression constructing the local derivation.
\end{itemize}
The elaborator performs only metalevel tasks: arity checking, topological sorting, variable allocation, common-subexpression sharing, generation of primitive-recursive symbols, and construction of a proof schedule.  It is acceptable for the elaborator to fail.  Its output, when accepted by Lean, is checked independently by the kernel.

\subsection{The two-pass compiler}

For a terminal query \(q_*\), compilation has two passes.

\paragraph{Reverse pass.}
Start with the terminal demand.  Traverse the acyclic graph backward and apply \cref{lem:demand-least}.  A loop node may contain an internal backward recurrence, but its horizon is bounded by a previously computed fuel term.  The result is a finite demand table \(D_G(x,q_*)\).

\paragraph{Forward pass.}
Traverse the graph in topological order.  Exact nodes compute terms.  Boolean and search nodes compute finite certificates.  Loop nodes run their specialized state recurrences.  Relation nodes compute viability tables for the finite horizons demanded by consumers.  Observer nodes extract only stable terminal information.  By \cref{thm:one-loop}, the entire pass may be represented by one packed state recurrence.

\paragraph{Proof pass.}
The base and step proofs establish the combined trace invariant.  At the final state, each node's exported formula is available.  The metalevel elaborator instantiates the local node rules and arranges them in a kernel proof schedule.  \Cref{thm:schedule} produces the terminal derivation.

\subsection{Soundness theorem}

We can now prove \cref{thm:main-compilation} precisely.

\begin{proof}[Proof of \cref{thm:main-compilation}]
Fix the graph, its parameters, and a terminal query.

\emph{Primitive-recursive specialization.}  Since the graph has finitely many nodes and every numerical constituent is primitive recursive, reverse demand is primitive recursive by reverse finite recursion.  Each forward node is primitive recursive: exact composition is closed under composition; Boolean folds, search, and viability are bounded recursions; each loop node is supplied as primitive recursion.  The one-loop construction of \cref{thm:one-loop} therefore yields a specialized primitive-recursive state symbol \(S_G\).

\emph{Quantifier-free invariant.}  Replace every bounded formula appearing in a local certificate by its Boolean code using \cref{thm:booleanization}.  Add range equations for every flag, index, and code.  Insert the supplied loop invariants and local transition equations, guarded by the finite program-counter cases.  The resulting formula \(I_G(t,S_G(t),x,q_*)\) is quantifier-free.

\emph{Trace proof.}  The initializer equations and local base certificates prove the base case.  For the step, use Boolean reflection to split on the program counter.  In each branch, invoke the corresponding local preservation derivation, rewrite by the transition equations, and rebuild the finite conjunction.  \Cref{thm:trace-reflection} gives the invariant at the terminal fuel.

\emph{Node certificates.}  Terminal projection equations extract the equations, Boolean facts, search facts, relational validity facts, and observer facts exported by each node.  Each such fact is appended to a kernel proof schedule.  Whenever a local rule needs predecessor facts, they are already earlier in the schedule by topological ordering.  Instantiate its variables by the generated state projections, then apply substitution.

\emph{Composition.}  Compile the schedule by \cref{thm:schedule}.  This yields \(\Gamma_G\vdash\Phi_G\).  Every derivation leaf is one of: a \(\PRAAx\) instance, a supplied local certificate, or a hypothesis from \(\Gamma_G\).  The only global principles used are intuitionistic natural deduction, equality substitution, simultaneous cut, and QF induction.  No arbitrary program code is interpreted and no unbounded witness is selected.
\end{proof}

\subsection{A relative completeness statement}

The compiler is complete for its declared node language in a modest but useful sense.

\begin{proposition}[Structural adequacy]
Let a finite proof construction be written using the six node classes, with explicit bounds, demand transformers, local derivations, and QF loop invariants.  Then there is a \KIFC graph whose terminal formula and derivation are definitionally the same construction up to state coding, substitution, and derived cut.
\end{proposition}

\begin{proof}
Introduce one node for each source combinator and preserve the source dependency order.  Replace each bounded control construct by its corresponding search, fold, relation, or loop node.  Replace nested control by the one-loop state encoding.  The local source derivations become node rules.  State coding changes terms but not the theorem: constructor/projection equations transport every local formula to the packed representation.  The proof schedule then reproduces the original sequence of local inferences.
\end{proof}

This is not a claim that every proof of mathematics belongs to the node language.  It says that once a proposed finitary proof has been reduced to these standard operations, the remaining passage to an embedded derivation is uniform.

\subsection{Why the compiler is not a universal evaluator}

It is tempting to define one object-level function \(U(e,x)\) that evaluates any encoded primitive-recursive program.  That cannot be total and primitive recursive while covering all unary primitive-recursive functions.

\begin{barrier}[No total universal primitive-recursive evaluator]\label{barrier:universal}
There is no primitive-recursive \(U(e,x)\) and coding of all unary primitive-recursive functions such that every unary primitive-recursive \(f\) has a code \(e\) with \(U(e,x)=f(x)\) for all \(x\).
\end{barrier}

\begin{proof}
If such \(U\) existed, the diagonal function \(g(e)=U(e,e)+1\) would be primitive recursive.  Let \(d\) be its code.  Then
\[
 g(d)=U(d,d)+1=g(d)+1,
\]
a contradiction.
\end{proof}

\KIFC avoids the premise.  Graph syntax is consumed by a metalevel elaborator, and the emitted object term is a fixed composition/recursion tree.  A clocked evaluator for one finite instruction set and one explicit step bound is also harmless, because it verifies only a bounded trace and does not enumerate all total primitive-recursive functions extensionally.

\subsection{Proof provenance}

A generated theorem should include machine-readable provenance:
\begin{itemize}
  \item repository commit and module imports;
  \item graph schema version and graph hash;
  \item list of local theorem names used by each node;
  \item generated formula name and public theorem name;
  \item \texttt{\#print axioms} output;
  \item clean-build result and absence of \texttt{sorry}.
\end{itemize}
This turns a large generated proof into an auditable dependency graph rather than an opaque tactic success.
